% Ad Atticum 16.2 --- headnote
% ad-atticum-16-02, 11 July 44 BC, in Puteolano

\textit{Cicero to Atticus, written at the Puteolan villa on
11 July 44 BC --- Perseus dateline \textit{Scr. in Puteolano
v Id. Quint. a. 710 (44)}, three days after 16.1. Cicero is
still at Puteoli, still preparing to embark, still being
hounded by creditors and family business while the political
situation reshapes itself daily. The first paragraphs are
financial: Hortensius (the orator's son and Cicero's debtor in
reverse direction here) is being shameless about a payment;
Publilius --- Cicero's former brother-in-law from the brief
marriage to Publilia --- is owed money out of the Tullia estate
settlement, and Cicero has already paid down 200{,}000
sesterces of a 400{,}000 outstanding to make the dissolution go
through. Atticus is being asked to govern Cicero's affairs at
Rome with full authority, even, if necessary, to liquidate
property to keep Cicero's good name intact.}

\textit{The political paragraphs surround Brutus's games. Brutus,
as praetor in absentia, has staged the Ludi Apollinares; the
play put on was Accius's \textit{Tereus}, with lines about
tyranny that the audience evidently applauded. Cicero is glad
of the gesture but bitterly remarks that the Roman people are
``spending their hands not in defending the commonwealth but in
applauding,'' and breaks into a verse cap: \textit{provided
they suffer something, let them suffer anything at all}. The
decision to sail is being praised by everyone, and Cicero ---
``thrust out with a pitchfork'' --- is now set on Brundisium
rather than Puteoli, having decided the legions are easier to
avoid than the pirates. Cassius has arrived at Naples with his
little squadron. The letter closes with the year's two
philosophical projects: the (now lost) \textit{Heracleides}
treatise to be tackled if Brundisium goes well, and the
\textit{De Gloria}, just sent off to Atticus with instructions
for Salvius to read excerpts at dinners --- only after a warm
audience has been secured.}
